\chapter*{Заключение}
\addcontentsline{toc}{chapter}{Заключение}
В ходе работы получены следующие результаты, соответствующие аналитической, теоретической, инженерной и практической составляющим исследования.\par
\begin{compactenum}
  \item Аналитические результаты состоят в систематизации задачи автоматизированного ранжирования кандидатов в контуре подбора персонала, анализе роли систем автоматизированного подбора персонала, сравнении лексических, семантических и нейросетевых подходов к сопоставлению вакансий и резюме, а также в выделении ограничений, связанных с качеством исторических кадровых данных, позиционным смещением и необходимостью контроля решений специалистом по подбору персонала.
  \item Теоретические результаты включают формализацию задачи как задачи информационного поиска и ранжирования: описание вакансии рассматривается как запрос, цифровые профили кандидатов и резюме рассматриваются как документы, а релевантность задается через историческое продвижение кандидата на следующий этап. Для оценки качества обоснованы метрики нормализованной дисконтированной кумулятивной выгоды, точности и полноты на первых позициях, а также средней обратной позиции первого релевантного результата. В качестве базового метода выбран подход раздельного кодирования вакансии и резюме (bi-encoder), при котором итоговая оценка рассчитывается по косинусному сходству плотных векторных представлений.
  \item Инженерные результаты представлены воспроизводимым программным контуром: подготовлены связанные обучающие примеры текст вакансии -- текст резюме с бинарными метками кадрового события, выполнена доменная настройка модели на базе \path|cointegrated/rubert-tiny2|, реализован модуль ATSRanker с жесткой фильтрацией, кодированием текстов через Sentence Transformer и сортировкой кандидатов по семантическому сходству, подготовлен сервис программного интерфейса приложения (Application Programming Interface, API) на FastAPI и выполнено подключение к тестовому контуру системы автоматизированного подбора персонала Reqcore.
  \item Практические результаты подтверждают применимость разработанного модуля как инструмента предварительной сортировки кандидатов без автономного принятия кадрового решения. На полном файле тестовой выборки \path|data/splits/test.tsv| для модели раздельного кодирования получены значения NDCG@3 = 1.0000 на 33 вакансиях, Precision@3 = 0.9958, Recall@3 = 0.9824 и MRR = 1.0000 на 119 вакансиях, где был хотя бы один кандидат, принятый к рассмотрению или приглашенный на следующий этап. Эти результаты показывают полезность доменной настройки для текущего корпуса, но должны интерпретироваться с учетом разреженности тестовой выборки и необходимости последующей промышленной проверки.
\end{compactenum}
Практическая ценность работы заключается в том, что разработанный модуль оформляет воспроизводимый путь от сырых текстов вакансий и резюме к числовой оценке семантического соответствия, которую можно встроить в существующий контур анализа заявки. При этом границы применимости зафиксированы явно: модуль должен использоваться как средство поддержки специалиста по подбору персонала с контролем качества, объяснимости и возможных смещений, а не как самостоятельный механизм принятия кадрового решения.\par
Направления будущей работы включают:\par
\begin{compactenum}
  \item Масштабирование поиска кандидатов через приближенный поиск ближайших соседей (Approximate Nearest Neighbor, ANN) и подключение векторной базы данных (Vector Database) для работы с большими индексами резюме.
  \item Развитие обучения ранжированию (Learning to Rank), включая попарные и списковые методы, например RankNet, LambdaRank или LambdaMART, а также расширенное сравнение с перекрестными моделями совместного кодирования.
  \item Добавление объяснений результата ранжирования для специалиста по подбору персонала: выделение совпавших навыков, опыта, требований вакансии и факторов, повлиявших на позицию кандидата.
  \item Проведение аудита справедливости и устойчивости ранжирования, включая анализ возможных смещений в исторических данных и проверку механизмов справедливого ранжирования результатов (fair ranking).
  \item Промышленное нагрузочное тестирование, мониторинг качества модели после внедрения и регламент периодической переоценки на новых данных.
\end{compactenum}
